\Needspace{8\baselineskip}
\section{Data and Evaluation Protocol}\label{sec:data}
\subsection{Existing Data and Conservative Independence}
The collection contains 418 photographs and 152 frames extracted from videos of three existing mid-category tyres, yielding 570 images. Acquisition records identify 12 physical tyres. The frames add views of existing subjects, and tyre identity relies on the acquisition record.

The training allocation contains 386 images, validation 81, and the historical test allocation 103 (Table~\ref{tab:data}). Images carry low-, mid-, or high-mileage proxy categories derived from source odometer folders. Low combines new tyres with the 100--5,000 group; mid contains the 40,000, 70,000, and 90,000 groups; high is the 100,000-plus group. These source labels are not measurements of abrasion. The known tyre allocation is eight/two/two. Unresolved relationships among three mid-category video parents are conservatively merged for component-based analysis. Consequently, the training role has six conservative connected components: one low, one mid, and four high. Their image counts are 67, 249, 60, 5, 3, and 2. This grouping does not assert that the three user-identified mid tyres are physically identical; it prevents uncertain provenance from being treated as independent evidence.

\begin{table}[t]
\caption{Fixed image allocation and supervision. Components are conservative provenance groups, not additional tyres.}\label{tab:data}
\centering\small
\setlength{\tabcolsep}{4pt}
\begin{tabularx}{\columnwidth}{Lrrr}\toprule
Quantity & Train & Validation & Historical test\\\midrule
Images & 386 & 81 & 103\\
Low / mid / high & 67 / 249 / 70 & 40 / 0 / 41 & 62 / 0 / 41\\
User-confirmed tyres & 8 & 2 & 2\\
Conservative components & 6 & 2 & 2\\
Human point images & 70 & 25 & 25\\
Weak point images & 152 & 0 & 0\\
No point target & 164 & 56 & 78\\\bottomrule
\end{tabularx}
\end{table}

Neither held-out role contains mid-category ground truth. Observed-class balanced accuracy (BA) is therefore the mean of low and high recalls, even though a classifier can predict all three categories. It is not three-class coverage. Class and component identity are strongly confounded: each held-out class corresponds to one conservative component per role. Group-aware evaluation is essential in dependent data \cite{groups}; here the small number of independent groups also limits what can be inferred after grouping.

\subsection{Geometry and Annotation Provenance}
The 418 original images have indexed masks, and the 152 frames have polygon-derived masks. Tyre and tread targets are overlapping binary channels. Pixels contradicted by the annotation contract are ignored rather than forced into mutually exclusive classes. A 120-image subset provides six human boundary points per image, allocated 70/25/25. The 152 video frames additionally supply weak polygon-derived point targets, weighted 0.25 during training. These weak targets are not counted as independent human evaluation truth.

Mean Dice describes agreement with available segmentation annotations; horizontal point mean absolute error (MAE) is normalized by image width and evaluated on human point images. Image-relative boundary positions do not provide physical camber, toe, or groove-depth measurements. TRACE uses frozen predicted support, not ground-truth masks at inference. Annotation ignore regions affect target-loss eligibility and do not enter the texture network as observation channels.

\subsection{Development History and Reporting Units}
The development sequence comprises ten training runs of 60 epochs each. TREAD-ONE has three seeds, the shared-gradient diagnostic has two matched runs, and ATLAS, ATLAS-LC, ATLAS-CONTEXT, GEO, and TRACE have one run each. The spatial audit uses TRACE at epoch 60.

Later candidates have one seed each. Their sequence is hypothesis-driven development rather than a factorial or fully seed-replicated comparison. Hyperparameters and selection contracts were frozen separately for each execution, but the overall sequence responded to previously inspected outcomes. Model-selection bias can arise even with a fixed finite validation allocation \cite{selectionbias}. We therefore distinguish validation-selected checkpoints from fixed final endpoints and explicitly acknowledge repeated historical-test inspection. No test-best checkpoint replaces the declared validation choice.

The primary new investigation is post-hoc analysis of retained TRACE evidence. Its nine interventions and arithmetic were recorded before that calculation; they were not prospectively preregistered before the model was trained. The mask-layout addendum was formulated later, after a high masked-token rank was observed, and is labeled exploratory. No image-independent confidence intervals or significance tests are reported: treating the many related images as independent replicates would misrepresent the support. Image means, component summaries, complete counts, and exact losses are reported instead.
